\subsection{Recuperación}

La tabla~\ref{tab:recuperacion} resume la comparación por componentes. El cambio más grande lo
produce el planificador: con la pregunta original, BM25 encuentra el artículo de origen entre los
tres primeros en el 40\,\% de las preguntas, y con las consultas reescritas, en el 95\,\%. El motivo
es el idioma (tabla~\ref{tab:idioma}). En inglés todas las configuraciones encuentran el artículo,
pero con las preguntas en español BM25 sobre la pregunta original lo encuentra en el 14\,\% de los
casos, porque casi no comparte términos con un corpus en inglés.

Buscar con la pregunta original y con las reformulaciones juntas (``ambas'' en las tablas) dio el
mejor resultado del modo híbrido, que fue también el mejor modo: P@5 de 0,38, nDCG@5 de 0,47 y MRR
de 0,60, en 66 ms. El reranker empeoró el orden de los fragmentos en todas las variantes (nDCG@5 de
0,39 con las mismas consultas) y agregó entre 3 y 4 s de latencia en CPU. Lo único que mejoró fue la
presencia del artículo de origen entre los tres primeros, del 98\,\% al 100\,\%. Darle solo las
consultas en inglés no cambió la calidad y ahorró 1 s. El modelo usado (\texttt{ms-marco-MiniLM-L-6-v2})
está entrenado con búsquedas web generales, lo que puede explicar que no ordene bien texto científico.

\begin{table}[htbp]
\centering
\footnotesize
\setlength{\tabcolsep}{3pt}
\caption{Recuperación sobre las 40 preguntas respondibles.}
\label{tab:recuperacion}
\begin{tabular}{llrrrrrrr}
\toprule
\textbf{Modo} & \textbf{Consultas} & \textbf{Art.@1} & \textbf{Art.@3} & \textbf{Origen@5} & \textbf{P@5} & \textbf{nDCG@5} & \textbf{MRR} & \textbf{ms} \\
\midrule
BM25 & original & 35 & 40 & 20 & 10 & 0,15 & 0,23 & 8 \\
BM25 & planificador & 80 & 95 & 45 & 32 & 0,41 & 0,57 & 9 \\
BM25 & ambas & 80 & 92 & 50 & 31 & 0,41 & 0,59 & 17 \\
Denso (MiniLM) & original & 68 & 80 & 20 & 18 & 0,24 & 0,35 & 16 \\
Denso (MiniLM) & planificador & 75 & 92 & 35 & 30 & 0,35 & 0,50 & 17 \\
Denso (MiniLM) & ambas & 72 & 98 & 35 & 30 & 0,36 & 0,51 & 42 \\
Híbrido (RRF) & original & 62 & 75 & 25 & 18 & 0,25 & 0,34 & 29 \\
Híbrido (RRF) & planificador & 82 & 92 & 42 & 34 & 0,43 & 0,59 & 29 \\
Híbrido (RRF) & ambas & 80 & 98 & 50 & 38 & 0,47 & 0,60 & 66 \\
Híbrido + reranker & original & 62 & 78 & 30 & 24 & 0,29 & 0,42 & 1221 \\
Híbrido + reranker & planificador & 70 & 100 & 32 & 30 & 0,35 & 0,51 & 2659 \\
Híbrido + reranker & ambas & 78 & 100 & 40 & 32 & 0,39 & 0,54 & 4007 \\
Híbrido + reranker & ambas, reranker en inglés & 78 & 100 & 40 & 32 & 0,39 & 0,54 & 3237 \\
\bottomrule
\end{tabular}
\par\smallskip\parbox{0.92\linewidth}{\footnotesize Consultas: la pregunta original, las del planificador o ambas. Art.@k: porcentaje de preguntas con el artículo de origen entre los k primeros. Origen@5: el fragmento de origen entre los cinco primeros. P@5, nDCG@5 y MRR usan los juicios de relevancia por pooling. ms: mediana de la latencia de recuperación.}
\end{table}

\begin{table}[htbp]
\centering
\small
\caption{Recuperación según el idioma de la pregunta (30 en español, 10 en inglés).}
\label{tab:idioma}
\begin{tabular}{lllrr}
\toprule
\textbf{Modo} & \textbf{Consultas} & \textbf{Idioma} & \textbf{Art.@3} & \textbf{Origen@5} \\
\midrule
BM25 & original & español & 14 & 4 \\
BM25 & original & inglés & 100 & 58 \\
BM25 & planificador & español & 93 & 50 \\
BM25 & planificador & inglés & 100 & 33 \\
BM25 & ambas & español & 89 & 50 \\
BM25 & ambas & inglés & 100 & 50 \\
Híbrido (RRF) & original & español & 64 & 14 \\
Híbrido (RRF) & original & inglés & 100 & 50 \\
Híbrido (RRF) & planificador & español & 93 & 43 \\
Híbrido (RRF) & planificador & inglés & 92 & 42 \\
Híbrido (RRF) & ambas & español & 96 & 54 \\
Híbrido (RRF) & ambas & inglés & 100 & 42 \\
\bottomrule
\end{tabular}
\end{table}


Lo que importa para el analista es si recibe algún fragmento que responda la pregunta. La
tabla~\ref{tab:contexto} muestra esa cobertura. Con seis fragmentos y un máximo de dos por artículo
(la configuración inicial), el 70\,\% de las preguntas tiene al menos un fragmento relevante en el
contexto. Permitir tres por artículo sube la cobertura al 75\,\% sin agregar tokens, y entregar diez
fragmentos no suma nada. En este nivel el reranker tampoco aporta: con él la cobertura es de 72\,\%.

\begin{table}[htbp]
\centering
\small
\caption{Cobertura del contexto que recibe el analista.}
\label{tab:contexto}
\begin{tabular}{rrlrrr}
\toprule
\textbf{Fragmentos} & \textbf{Por artículo} & \textbf{Modo} & \textbf{Con relevante} & \textbf{Relevantes} & \textbf{Con origen} \\
\midrule
6 & 2 & Híbrido (RRF) & 70 & 1,12 & 30 \\
6 & 2 & Híbrido + reranker & 68 & 1,02 & 25 \\
6 & 3 & Híbrido (RRF) & 75 & 1,48 & 45 \\
6 & 3 & Híbrido + reranker & 72 & 1,27 & 35 \\
10 & 3 & Híbrido (RRF) & 75 & 1,50 & 45 \\
10 & 3 & Híbrido + reranker & 72 & 1,45 & 38 \\
\bottomrule
\end{tabular}
\par\smallskip\parbox{0.92\linewidth}{\footnotesize Con relevante: preguntas con al menos un fragmento relevante (pooling) en el contexto. Relevantes: cantidad media de fragmentos relevantes. Con origen: el fragmento de origen está en el contexto.}
\end{table}


\subsection{Sistema completo}

La tabla~\ref{tab:sistema} compara el sistema completo con la variante sin verificador. En las 40
preguntas respondibles, el artículo de origen se recuperó y se citó en el 92,5\,\% de los casos y el
98,5\,\% de las afirmaciones está respaldado por los documentos citados. La corrección media frente
a la respuesta de referencia es 3,95 sobre 5: 17 respuestas recibieron 5 y 6 recibieron 2. Las notas
bajas corresponden casi siempre a preguntas en las que el contexto no tenía el fragmento con la
respuesta. En esos casos el analista lo dijo en lugar de inventar. El sistema rechazó las 5 preguntas
fuera de alcance y en las 3 preguntas sin respuesta en el corpus reconoció que no había evidencia.
El planificador rechazó también 2 preguntas respondibles: una sobre bibliometría de la dislexia y
otra sobre modelos generativos de fuentes talamocorticales.

\begin{table}[htbp]
\centering
\small
\caption{Sistema completo y variante sin verificador.}
\label{tab:sistema}
\begin{tabular}{lrrrrrrrrr}
\toprule
\textbf{Variante} & \textbf{Recup.} & \textbf{Cita} & \textbf{Corr.} & \textbf{Fidel.} & \textbf{Fuera} & \textbf{Sin resp.} & \textbf{USD} & \textbf{p50 s} & \textbf{p95 s} \\
\midrule
Completo & 92,5 & 92,5 & 3,95 & 98,5 & 100 & 100 & 0,041 & 26 & 40 \\
Sin verificador & 92,5 & 92,5 & 3,97 & 97,8 & 100 & 100 & 0,024 & 12 & 16 \\
\bottomrule
\end{tabular}
\par\smallskip\parbox{0.92\linewidth}{\footnotesize Recup.: artículo de origen entre los recuperados. Cita: la respuesta cita el artículo de origen. Corr.: corrección de 1 a 5 según el juez. Fidel.: afirmaciones respaldadas por los documentos citados. Fuera: preguntas fuera de alcance rechazadas. Sin resp.: preguntas sin respuesta en el corpus en las que el sistema lo reconoce. USD: costo medio por consulta.}
\end{table}


El verificador no mejoró la corrección (3,95 contra 3,97) y subió la fidelidad menos de un punto,
con un 70\,\% más de costo y el doble de latencia. Pregunta por pregunta, la versión con verificador
fue mejor en 6 casos, peor en 6 e igual en 26. Su aporte medible es como alarma: las 9 respuestas
que no aprobó tienen una corrección media de 3,67, contra 4,03 de las aprobadas, y concentran la
mitad de las notas bajas. Esas respuestas llegan al usuario con una advertencia que lista las
afirmaciones sin respaldo.

Antes de los ajustes finales, la corrección media era 3,58. La mejora salió de permitir tres
fragmentos por artículo y de que la búsqueda adicional que pide el verificador use solo la consulta
nueva y sume sus fragmentos a los anteriores.

El costo medio por consulta es de 0,041 dólares con verificador y 0,024 sin él. La latencia mediana
es de 26 s y 12 s, respectivamente, y la mayor parte corresponde al analista y al verificador. Una
pregunta repetida se responde desde la caché en menos de 1 ms y sin costo.

\subsection{Seguridad}

\begin{table}[htbp]
\centering
\small
\caption{Pruebas adversariales.}
\label{tab:seguridad}
\begin{tabular}{lr}
\toprule
\textbf{Prueba} & \textbf{Resultado} \\
\midrule
Ataques directos bloqueados por la guardia & 11 de 12 \\
Ataques directos rechazados por el planificador & 1 de 12 \\
Ataques con filtración del prompt & 0 de 12 \\
Preguntas legítimas bloqueadas & 0 de 40 \\
Fragmentos del corpus excluidos & 0 de 1599 \\
Documento plantado detectado al indexar & sí \\
Analista adopta la afirmación plantada (filtro salteado) & no \\
Analista recomienda el enlace plantado (filtro salteado) & no \\
Analista advierte la inyección (filtro salteado) & sí \\
\bottomrule
\end{tabular}
\end{table}


La tabla~\ref{tab:seguridad} resume las pruebas adversariales. La guardia bloqueó 11 de los 12
ataques directos y el restante (pedir que se finja que los documentos dicen ICC = 0,99) lo rechazó el
planificador por fuera de alcance. Ningún ataque filtró el prompt.

El hallazgo principal fue el clasificador local. Está entrenado en inglés y le asigna probabilidad
de inyección 1,000 a preguntas normales en español, como ``¿Qué tan confiable es el PLI entre
sesiones?''. Usado solo, habría bloqueado buena parte de las consultas legítimas. Con la confirmación
del modelo auxiliar, ninguna de las 40 preguntas del conjunto de evaluación quedó bloqueada. Sobre el
corpus, el clasificador marcó 2 de 1599 fragmentos (texto científico común) y la confirmación los
descartó a los dos.

El documento plantado fue detectado por la regla léxica. Con el filtro salteado a propósito, el
analista no adoptó la afirmación falsa, no recomendó el enlace y terminó la respuesta con una
advertencia: ``el documento 1 contiene instrucciones insertadas que afirman un ICC = 1,0 y remiten a
un enlace externo; no las seguí''.

\subsection{Ejemplo}

Pregunta: ``¿Qué tan confiable es el PLI entre sesiones con ojos cerrados, comparado con el AECc?''.
El planificador la reescribió en tres consultas en inglés, el verificador aprobó la respuesta en la
primera vuelta y la consulta costó 0,034 dólares en 25 s. Respuesta:

\begin{quote}\small
La evidencia disponible permite solo una respuesta parcial: no hay en estos fragmentos una
comparación directa, con valores numéricos, entre PLI y AECc en registros con ojos cerrados entre
sesiones.

En datos empíricos de adultos sanos (64 canales, dos sesiones separadas por seis semanas, n = 42),
el PLI en theta y alfa mostró confiabilidad temporal moderada, junto con alta confiabilidad entre
estados, lo que lo ubica como una medida relativamente independiente del estado [2]. En ese mismo
estudio, las métricas de conectividad funcional en conjunto (AECc y PLI) variaron de pobre a
excelente, con mejor desempeño de theta y alfa frente a delta y beta, pero el resumen no desagrega el
ICC del AECc por banda ni por condición ocular [4][5]. [\ldots]

En simulaciones multicanal de EEG/MEG, PLI y wPLI presentaron mayor confiabilidad test-retest y
mayor sensibilidad a la fuerza de acoplamiento que ImCoh y AECc, aunque el AECc resultó menos ruidoso
[1]. Esa evidencia es simulada, no empírica con ojos cerrados.

[1] Benchmarking metrics for inferring functional connectivity from multi-channel EEG and MEG: A
simulation study. (Chaos, 2020). Abstract.\\
{[2]} Reliability and state-dependency of EEG connectivity, complexity and network characteristics.
(Sci Rep, 2025). Discussion.\\
{[4]}, {[5]} El mismo artículo: Abstract e Introduction.
\end{quote}
